\input{preamble}

% OK, start here.
%
\begin{document}

\title{Limites de champs algébriques}


\maketitle

\phantomsection
\label{section-phantom}

\tableofcontents

\section{Introduction}
\label{section-introduction}

\noindent
Dans ce chapitre, nous rassemblons des résultats relatifs aux limites de
champs algébriques. De nombreux résultats sur les limites de champs
algébriques et d'espaces algébriques ont été obtenus par David Rydh dans
\cite{rydh_approx}.


\section{Conventions}
\label{section-conventions}

\noindent
Nous continuons d'employer les conventions et l'abus de langage introduits
dans Propriétés des champs, section
\ref{stacks-properties-section-conventions}.







\section{Morphismes de présentation finie}
\label{section-finite-presentation}

\noindent
Cette section est l'analogue de Limites d'espaces, section
\ref{spaces-limits-section-finite-presentation}. Nous y avons défini ce que
signifie, pour une transformation de foncteurs sur $\Sch$, le fait de commuter
aux limites (nous suggérons de consulter la caractérisation de Limites
d'espaces, lemme
\ref{spaces-limits-lemma-characterize-relative-limit-preserving}). Dans
Critères de représentabilité, section \ref{criteria-section-limit-preserving},
nous avons défini la notion de « commutation aux limites sur les objets ».
Rappelons enfin que, dans Axiomes d'Artin, section \ref{artin-section-limits},
nous avons défini ce que signifie, pour une catégorie fibrée en groupoïdes sur
$\Sch$, le fait de commuter aux limites. En combinant ces notions, nous
obtenons la définition suivante.

\begin{definition}
\label{definition-limit-preserving}
Soit $S$ un schéma. Soit $f : \mathcal{X} \to \mathcal{Y}$ un
$1$-morphisme de catégories fibrées en groupoïdes sur $(\Sch/S)_{fppf}$.
Nous disons que $f$ {\it commute aux limites} si, pour toute limite projective
filtrante $U = \lim U_i$ de schémas affines sur $S$, le diagramme
$$
\xymatrix{
\colim \mathcal{X}_{U_i} \ar[r] \ar[d]_f & \mathcal{X}_U \ar[d]^f \\
\colim \mathcal{Y}_{U_i} \ar[r] & \mathcal{Y}_U
}
$$
de catégories fibres est $2$-cartésien.
\end{definition}

\begin{lemma}
\label{lemma-limit-preserving-objects}
Soit $S$ un schéma. Soit $f : \mathcal{X} \to \mathcal{Y}$ un
$1$-morphisme de catégories fibrées en groupoïdes sur $(\Sch/S)_{fppf}$.
Si $f$ commute aux limites (définition \ref{definition-limit-preserving}),
alors $f$ commute aux limites sur les objets (Critères de représentabilité,
section \ref{criteria-section-limit-preserving}).
\end{lemma}

\begin{proof}
Si, pour toute limite projective filtrante $U = \lim U_i$ de schémas affines
sur $U$, le foncteur
$$
\colim \mathcal{X}_{U_i} \longrightarrow
(\colim \mathcal{Y}_{U_i}) \times_{\mathcal{Y}_U} \mathcal{X}_U
$$
est essentiellement surjectif, alors $f$ commute aux limites sur les objets.
\end{proof}

\begin{lemma}
\label{lemma-base-change-limit-preserving}
Soient $p : \mathcal{X} \to \mathcal{Y}$ et
$q : \mathcal{Z} \to \mathcal{Y}$ des $1$-morphismes de catégories fibrées
en groupoïdes sur $(\Sch/S)_{fppf}$. Si
$p : \mathcal{X} \to \mathcal{Y}$ commute aux limites, il en va de même du
changement de base
$p' : \mathcal{X} \times_\mathcal{Y} \mathcal{Z} \to \mathcal{Z}$
de $p$ suivant $q$.
\end{lemma}

\begin{proof}
C'est formel. Soit $U = \lim_{i \in I} U_i$ la limite projective filtrante
de schémas affines $U_i$ sur $S$. Pour tout $i$, nous avons
$$
(\mathcal{X} \times_\mathcal{Y} \mathcal{Z})_{U_i} =
\mathcal{X}_{U_i} \times_{\mathcal{Y}_{U_i}} \mathcal{Z}_{U_i}
$$
Les colimites filtrantes commutent aux $2$-produits fibrés de catégories
(nous omettons les détails). Puisque $p$ commute aux limites, il vient
\begin{align*}
\colim (\mathcal{X} \times_\mathcal{Y} \mathcal{Z})_{U_i}
& =
\colim \mathcal{X}_{U_i} \times_{\colim \mathcal{Y}_{U_i}}
\colim \mathcal{Z}_{U_i} \\
& =
\mathcal{X}_U \times_{\mathcal{Y}_U} \colim \mathcal{Y}_{U_i}
\times_{\colim \mathcal{Y}_{U_i}}
\colim \mathcal{Z}_{U_i} \\
& =
\mathcal{X}_U \times_{\mathcal{Y}_U} \colim \mathcal{Z}_{U_i} \\
& =
\mathcal{X}_U \times_{\mathcal{Y}_U} \mathcal{Z}_U \times_{\mathcal{Z}_U}
\colim \mathcal{Z}_{U_i} \\
& =
(\mathcal{X} \times_\mathcal{Y} \mathcal{Z})_U \times_{\mathcal{Z}_U}
\colim \mathcal{Z}_{U_i}
\end{align*}
comme voulu.
\end{proof}

\begin{lemma}
\label{lemma-composition-limit-preserving}
Soient $p : \mathcal{X} \to \mathcal{Y}$ et
$q : \mathcal{Y} \to \mathcal{Z}$ des $1$-morphismes de catégories fibrées
en groupoïdes sur $(\Sch/S)_{fppf}$. Si $p$ et $q$ commutent aux limites,
leur composé $q \circ p$ y commute également.
\end{lemma}

\begin{proof}
C'est formel. Soit $U = \lim_{i \in I} U_i$ la limite projective filtrante
de schémas affines $U_i$ sur $S$. Puisque $p$ et $q$ commutent aux limites,
nous obtenons
\begin{align*}
\colim \mathcal{X}_{U_i}
& =
\mathcal{X}_U \times_{\mathcal{Y}_U} \colim \mathcal{Y}_{U_i} \\
& =
\mathcal{X}_U \times_{\mathcal{Y}_U} \mathcal{Y}_U
\times_{\mathcal{Z}_U} \colim \mathcal{Z}_{U_i} \\
& =
\mathcal{X}_U \times_{\mathcal{Z}_U} \colim \mathcal{Z}_{U_i}
\end{align*}
comme voulu.
\end{proof}

\begin{lemma}
\label{lemma-representable-by-spaces-limit-preserving}
Soit $p : \mathcal{X} \to \mathcal{Y}$ un $1$-morphisme de catégories
fibrées en groupoïdes sur $(\Sch/S)_{fppf}$. Si $p$ est représentable par
des espaces algébriques, les conditions suivantes sont équivalentes :
\begin{enumerate}
\item $p$ commute aux limites,
\item $p$ commute aux limites sur les objets, et
\item $p$ est localement de présentation finie (voir Champs algébriques,
définition \ref{algebraic-definition-relative-representable-property}).
\end{enumerate}
\end{lemma}

\begin{proof}
Dans Critères de représentabilité, lemme
\ref{criteria-lemma-representable-by-spaces-limit-preserving}, nous avons vu
que (2) et (3) sont équivalentes. Il suffit donc de montrer l'équivalence de
(1) et (2). Une implication a été établie dans le lemme
\ref{lemma-limit-preserving-objects}. Pour la réciproque, soit
$U = \lim_{i \in I} U_i$ la limite projective filtrante de schémas affines
$U_i$ sur $S$. Nous devons montrer que
$$
\colim \mathcal{X}_{U_i} \longrightarrow
\mathcal{X}_U \times_{\mathcal{Y}_U} \colim \mathcal{Y}_{U_i}
$$
est une équivalence. L'hypothèse (2) assure qu'il est essentiellement
surjectif ; il reste donc à prouver qu'il est pleinement fidèle. Comme $p$
est fidèle sur les catégories fibres (Champs algébriques, lemme
\ref{algebraic-lemma-criterion-map-representable-spaces-fibred-in-groupoids}),
ce foncteur est fidèle. Soient $x_i$ et $x'_i$ des objets de la catégorie
fibre de $\mathcal{X}$ sur $U_i$. Le foncteur ci-dessus envoie $x_i$ sur
$(x_i|_U, p(x_i), can)$, où $can$ est l'isomorphisme canonique
$p(x_i|_U) \to p(x_i)|_U$. Supposons donc donné un morphisme
$$
(\alpha, \beta_i) : (x_i|_U, p(x_i), can) \longrightarrow
(x'_i|_U, p(x'_i), can)
$$
dans la catégorie qui figure au membre de droite de la première flèche
affichée de cette démonstration. Il s'agit de construire un $i' \geq i$ et un
morphisme $x_i|_{U_{i'}} \to x'_i|_{U_{i'}}$ dont l'image soit
$(\alpha, \beta_i|_{U_{i'}})$.

\medskip\noindent
Posons $y_i = p(x_i)$ et $y'_i = p(x'_i)$. D'après Champs algébriques,
lemme
\ref{algebraic-lemma-criterion-map-representable-spaces-fibred-in-groupoids},
le foncteur
$$
X_{y_i} : (\Sch/U_i)^{opp} \to \textit{Ensembles},\quad
V/U_i \mapsto
\{(x, \phi) \mid x \in \Ob(\mathcal{X}_V), \phi : f(x) \to y_i|V\}/\cong
$$
est un espace algébrique sur $U_i$, et il en va de même du foncteur
$X_{y'_i}$ défini de façon analogue. Puisque (2) équivaut à (3), nous voyons
que $X_{y'_i}$ est localement de présentation finie sur $U_i$. Remarquons que
$(x_i, \text{id})$ et $(x'_i, \text{id})$ définissent des points à valeurs
dans $U_i$ de $X_{y_i}$ et de $X_{y'_i}$. Il existe une transformation de
foncteurs
$$
\beta_i : X_{y_i} \to X_{y'_i},\quad
(x/V, \phi) \mapsto (x/V, \beta_i|_V \circ \phi)
$$
autrement dit, un morphisme d'espaces algébriques sur $U_i$. Nous affirmons
que le diagramme
$$
\xymatrix{
U \ar[d] \ar[rr] & & U_i \ar[d]^{(x'_i, \text{id})} \\
U_i \ar[r]^{(x_i, \text{id})} & X_{y_i} \ar[r]^{\beta_i} & X_{y'_i}
}
$$
est commutatif. En effet, cela équivaut à dire que les couples
$(x_i|_U, \beta_i|_U)$ et $(x'_i|_U, \text{id})$ intervenant dans la
définition du foncteur $X_{y'_i}$ sont isomorphes. Or le morphisme
$\alpha : x_i|_U \to x'_i|_U$ fournit précisément un tel isomorphisme. En
remontant cet argument, on voit qu'il suffit de trouver un $i' \geq i$ pour
lequel le diagramme
$$
\xymatrix{
U_{i'} \ar[d] \ar[rr] & & U_i \ar[d]^{(x'_i, \text{id})} \\
U_i \ar[r]^{(x_i, \text{id})} & X_{y_i} \ar[r]^{\beta_i} & X_{y'_i}
}
$$
soit commutatif : on obtient alors un isomorphisme
$x_i|_{U_{i'}} \to x'_i|_{U_{i'}}$ qui résout le problème posé au paragraphe
précédent. Or le morphisme diagonal
$$
\Delta : X_{y'_i} \to X_{y'_i} \times_{U_i} X_{y'_i}
$$
est localement de présentation finie (Morphismes d'espaces, lemme
\ref{spaces-morphisms-lemma-diagonal-morphism-finite-type}). Puisque
$U \to U_i$ égalise les deux morphismes vers $X_{y'_i}$, il existe donc un
$i' \geq i$ tel que $U_{i'} \to U_i$ les égalise également ; voir Limites
d'espaces, proposition
\ref{spaces-limits-proposition-characterize-locally-finite-presentation}.
\end{proof}

\begin{lemma}
\label{lemma-limit-preserving-diagonal}
Soit $p : \mathcal{X} \to \mathcal{Y}$ un $1$-morphisme de catégories
fibrées en groupoïdes sur $(\Sch/S)_{fppf}$. Les conditions suivantes sont
équivalentes :
\begin{enumerate}
\item la diagonale
$\Delta : \mathcal{X} \to \mathcal{X} \times_\mathcal{Y} \mathcal{X}$
commute aux limites, et
\item pour toute limite projective filtrante $U = \lim U_i$ de schémas
affines sur $S$, le foncteur
$$
\colim \mathcal{X}_{U_i} \longrightarrow
\mathcal{X}_U \times_{\mathcal{Y}_U} \colim \mathcal{Y}_{U_i}
$$
est pleinement fidèle.
\end{enumerate}
En particulier, si $p$ commute aux limites, il en va de même de $\Delta$.
\end{lemma}

\begin{proof}
Soit $U = \lim U_i$ une limite projective filtrante de schémas affines sur
$S$. Nous affirmons que le foncteur
$$
\colim \mathcal{X}_{U_i} \longrightarrow
\mathcal{X}_U \times_{\mathcal{Y}_U} \colim \mathcal{Y}_{U_i}
$$
est pleinement fidèle si et seulement si le foncteur
$$
\colim \mathcal{X}_{U_i} \longrightarrow
\mathcal{X}_U \times_{(\mathcal{X} \times_\mathcal{Y} \mathcal{X})_U}
\colim (\mathcal{X} \times_\mathcal{Y} \mathcal{X})_{U_i}
$$
est une équivalence. Cela démontrera le lemme. Comme
$(\mathcal{X} \times_\mathcal{Y} \mathcal{X})_U =
\mathcal{X}_U \times_{\mathcal{Y}_U} \mathcal{X}_U$
et
$(\mathcal{X} \times_\mathcal{Y} \mathcal{X})_{U_i} =
\mathcal{X}_{U_i} \times_{\mathcal{Y}_{U_i}} \mathcal{X}_{U_i}$
il s'agit d'une assertion purement catégorique, que nous examinons au
paragraphe suivant.

\medskip\noindent
Soit $\mathcal{I}$ une catégorie filtrante d'indices. Soient
$(\mathcal{C}_i)$ et $(\mathcal{D}_i)$ des systèmes de groupoïdes sur
$\mathcal{I}$. Soit
$p : (\mathcal{C}_i) \to (\mathcal{D}_i)$ un morphisme de systèmes de
groupoïdes sur $\mathcal{I}$. Supposons donnés un foncteur de groupoïdes
$p : \mathcal{C} \to \mathcal{D}$ et des foncteurs
$f : \colim \mathcal{C}_i \to \mathcal{C}$ et
$g : \colim \mathcal{D}_i \to \mathcal{D}$
qui s'insèrent dans un diagramme commutatif
$$
\xymatrix{
\colim \mathcal{C}_i \ar[d]_p \ar[r]_f & \mathcal{C} \ar[d]^p \\
\colim \mathcal{D}_i \ar[r]^g & \mathcal{D}
}
$$
Nous affirmons alors que
$$
A : \colim \mathcal{C}_i \longrightarrow
\mathcal{C} \times_\mathcal{D} \colim \mathcal{D}_i
$$
est pleinement fidèle si et seulement si le foncteur
$$
B : \colim \mathcal{C}_i \longrightarrow
\mathcal{C}
\times_{\Delta, \mathcal{C} \times_\mathcal{D} \mathcal{C}, f \times_g f}
\colim (\mathcal{C}_i \times_{\mathcal{D}_i} \mathcal{C}_i)
$$
est une équivalence. Posons $\mathcal{C}' = \colim \mathcal{C}_i$ et
$\mathcal{D}' = \colim \mathcal{D}_i$. Puisque les $2$-produits fibrés
commutent aux colimites filtrantes, $A$ et $B$ deviennent les foncteurs
$$
A' : \mathcal{C}' \to \mathcal{C} \times_\mathcal{D} \mathcal{D}'
\quad\text{et}\quad
B' : \mathcal{C}' \longrightarrow
\mathcal{C}
\times_{\Delta, \mathcal{C} \times_\mathcal{D} \mathcal{C}, f \times_g f}
(\mathcal{C}' \times_{\mathcal{D}'} \mathcal{C}')
$$
Il suffit donc de prouver que, si
$$
\xymatrix{
\mathcal{C}' \ar[d]_p \ar[r]_f & \mathcal{C} \ar[d]^p \\
\mathcal{D}' \ar[r]^g & \mathcal{D}
}
$$
est un diagramme commutatif de groupoïdes, alors $A'$ est pleinement fidèle
si et seulement si $B'$ est une équivalence. Cela résulte de Catégories,
lemme \ref{categories-lemma-fully-faithful-diagonal-equivalence} (avec une
catégorie de base triviale, c'est-à-dire ponctuelle), car
$$
\mathcal{C}
\times_{\Delta, \mathcal{C} \times_\mathcal{D} \mathcal{C}, f \times_g f}
(\mathcal{C}' \times_{\mathcal{D}'} \mathcal{C}') =
\mathcal{C}'
\times_{A', \mathcal{C} \times_\mathcal{D} \mathcal{D}', A'}
\mathcal{C}'
$$
La démonstration est achevée.
\end{proof}

\begin{lemma}
\label{lemma-locally-finite-presentation-limit-preserving}
Soit $S$ un schéma. Soit $\mathcal{X}$ un champ algébrique sur $S$. Si
$\mathcal{X} \to S$ est localement de présentation finie, alors
$\mathcal{X}$ commute aux limites au sens de Axiomes d'Artin, définition
\ref{artin-definition-limit-preserving} (de manière équivalente, le morphisme
$\mathcal{X} \to S$ commute aux limites).
\end{lemma}

\begin{proof}
Choisissons un morphisme lisse surjectif $U \to \mathcal{X}$, où $U$ est un
schéma. Alors $U \to S$ est localement de présentation finie ; voir Morphismes
de champs, section \ref{stacks-morphisms-section-finite-presentation}. Nous
pouvons écrire $\mathcal{X} = [U/R]$ pour un groupoïde lisse en espaces
algébriques $(U, R, s, t, c)$ ; voir Champs algébriques, lemme
\ref{algebraic-lemma-stack-presentation}. Puisque $U$ est localement de présentation
finie sur $S$, l'espace algébrique $R$ l'est lui aussi sur $S$. Rappelons que
$[U/R]$ est le champ en groupoïdes sur $(\Sch/S)_{fppf}$ obtenu en passant au
champ associé à la catégorie fibrée en groupoïdes dont la catégorie fibre sur
$T$ est le groupoïde $(U(T), R(T), s, t, c)$. Comme les foncteurs $U$ et $R$
commutent aux limites (Limites d'espaces, proposition
\ref{spaces-limits-proposition-characterize-locally-finite-presentation}),
cette catégorie fibrée en groupoïdes commute aux limites. Il suffit donc de
montrer que le passage au champ fppf associé conserve cette propriété. C'est
bien le cas (indication : utiliser Topologies, lemme
\ref{topologies-lemma-limit-fppf-topology}). Nous donnons néanmoins ci-dessous
une démonstration directe, en utilisant que nous connaissons ici explicitement
le passage au champ associé.

\medskip\noindent
Soit $T = \lim T_\lambda$ une limite projective filtrante de schémas affines
sur $S$. Nous devons montrer que le foncteur
$$
\colim [U/R]_{T_\lambda} \longrightarrow [U/R]_T
$$
est une équivalence de catégories. Montrons d'abord qu'il est essentiellement
surjectif. Soit $x \in \Ob([U/R]_T)$. Groupoïdes en espaces, lemme
\ref{spaces-groupoids-lemma-quotient-stack-objects}, décrit la catégorie
$[U/R]_T$. En particulier, $x$ correspond à un recouvrement fppf
$\{T_i \to T\}_{i \in I}$ et à une donnée de descente dans $[U/R]$
$(u_i, r_{ij})$ relative à ce recouvrement. Après avoir raffiné celui-ci, nous
pouvons supposer qu'il s'agit d'un recouvrement fppf standard du schéma affine
$T$. Par Topologies, lemme \ref{topologies-lemma-limit-fppf-topology}, nous
pouvons choisir un $\lambda$ et un recouvrement fppf standard
$\{T_{\lambda, i} \to T_\lambda\}_{i \in I}$ dont le changement de base à
$T$ soit égal à $\{T_i \to T\}_{i \in I}$. Pour tout $i$, quitte à augmenter
$\lambda$, il existe un morphisme $u_{\lambda, i} : T_{\lambda, i} \to U$
dont le composé avec $T_i \to T_{\lambda, i}$ est le morphisme donné $u_i$
(c'est ici que nous utilisons la commutation de $U$ aux limites). De même,
pour tous $i, j$, quitte à augmenter $\lambda$, il existe un morphisme
$r_{\lambda, ij} : T_{\lambda, i} \times_{T_\lambda} T_{\lambda, j} \to R$
dont le composé avec $T_{ij} \to T_{\lambda, ij}$ est le morphisme donné
$r_{ij}$ (c'est ici que nous utilisons la commutation de $R$ aux limites).
Quitte à augmenter encore $\lambda$, nous pouvons supposer que
$$
s \circ r_{\lambda, ij} = u_{\lambda, i} \circ \text{pr}_0
\quad\text{et}\quad
t \circ r_{\lambda, ij} = u_{\lambda, j} \circ \text{pr}_1,
$$
et
$$
c \circ (r_{\lambda, jk} \circ \text{pr}_{12}, r_{\lambda, ij}
\circ \text{pr}_{01}) = r_{\lambda, ik} \circ \text{pr}_{02}.
$$
Autrement dit, nous pouvons supposer que
$(u_{\lambda, i}, r_{\lambda, ij})$ est une donnée de descente dans $[U/R]$
relative au recouvrement
$\{T_{\lambda, i} \to T_\lambda\}_{i \in I}$. Nous obtenons alors un objet
correspondant de $[U/R]$ sur $T_\lambda$ dont l'image réciproque sur $T$ est
isomorphe à $x$, comme voulu. La pleine fidélité se démontre exactement de la
même façon, à l'aide de la description des morphismes dans les catégories
fibres de $[U/T]$ donnée dans Groupoïdes en espaces, lemme
\ref{spaces-groupoids-lemma-quotient-stack-objects}.
\end{proof}

\begin{proposition}
\label{proposition-characterize-locally-finite-presentation}
\begin{reference}
C'est un cas particulier de \cite[Lemme 2.3.15]{Emerton-Gee}.
\end{reference}
Soit $f : \mathcal{X} \to \mathcal{Y}$ un morphisme de champs algébriques.
Les conditions suivantes sont équivalentes :
\begin{enumerate}
\item $f$ commute aux limites,
\item $f$ commute aux limites sur les objets, et
\item $f$ est localement de présentation finie.
\end{enumerate}
\end{proposition}

\begin{proof}
Supposons (3). Soit $T = \lim T_i$ une limite projective filtrante de schémas
affines. Considérons le foncteur
$$
\colim \mathcal{X}_{T_i} \longrightarrow
\mathcal{X}_T \times_{\mathcal{Y}_T} \colim \mathcal{Y}_{T_i}
$$
Soit $(x, y_i, \beta)$ un objet du membre de droite, c'est-à-dire que
$x \in \Ob(\mathcal{X}_T)$, $y_i \in \Ob(\mathcal{Y}_{T_i})$ et que
$\beta : f(x) \to y_i|_T$ est un morphisme de $\mathcal{Y}_T$. Nous pouvons
alors considérer $(x, y_i, \beta)$ comme un objet du champ algébrique
$\mathcal{X}_{y_i} = \mathcal{X} \times_{\mathcal{Y}, y_i} T_i$ au-dessus de
$T$. Comme
$\mathcal{X}_{y_i} \to T_i$ est localement de présentation finie (c'est un
changement de base de $f$), il commute aux limites par le lemme
\ref{lemma-locally-finite-presentation-limit-preserving}. Ainsi,
$(x, y_i, \beta)$ provient d'un objet sur $T_{i'}$ pour un certain
$i' \geq i$. En explicitant les définitions, nous voyons que
$(x, y_i, \beta)$ appartient à l'image essentielle du foncteur affiché.
Celui-ci est donc essentiellement surjectif. Autrement dit, $f$ commute aux
limites sur les objets. Appliquons maintenant ce résultat à la diagonale
$\Delta$ de $f$. D'après Morphismes de champs, lemme
\ref{stacks-morphisms-lemma-diagonal-morphism-finite-type}, le morphisme
$\Delta$ est localement de présentation finie. L'argument précédent montre
donc que $\Delta$ commute aux limites sur les objets. Par le lemme
\ref{lemma-representable-by-spaces-limit-preserving}, $\Delta$ commute aux limites.
Le lemme \ref{lemma-limit-preserving-diagonal} entraîne alors que le foncteur
affiché ci-dessus est pleinement fidèle. C'est donc une équivalence, puisque
nous avons déjà prouvé sa surjectivité essentielle, et (1) est établi.

\medskip\noindent
L'implication (1) $\Rightarrow$ (2) est immédiate. Supposons (2). Choisissons
un schéma $V$ et un morphisme lisse surjectif $V \to \mathcal{Y}$. Par
Critères de représentabilité, lemme
\ref{criteria-lemma-base-change-limit-preserving}, le changement de base
$\mathcal{X} \times_\mathcal{Y} V \to V$ commute aux limites sur les objets.
Choisissons un schéma $U$ et un morphisme lisse surjectif
$U \to \mathcal{X} \times_\mathcal{Y} V$. Un morphisme lisse étant localement
de présentation finie, $U \to \mathcal{X} \times_\mathcal{Y} V$ commute aux
limites, d'après la première partie de la démonstration. Par Critères de
représentabilité, lemme \ref{criteria-lemma-composition-limit-preserving}, le
composé $U \to V$ commute aux limites sur les objets. Nous en concluons que
$U \to V$ est localement de présentation finie ; voir Critères de
représentabilité, lemme
\ref{criteria-lemma-representable-by-spaces-limit-preserving}. C'est exactement
la condition disant que $f$ est localement de présentation finie ; voir
Morphismes de champs, définition
\ref{stacks-morphisms-definition-locally-finite-presentation}.
\end{proof}




\section{Descente de propriétés}
\label{section-descent}

\noindent
Cette section est l'analogue de Limites, section \ref{limits-section-descent}.

\begin{situation}
\label{situation-descent}
Soit $Y = \lim_{i \in I} Y_i$ la limite projective d'un système filtrant
d'espaces algébriques dont les morphismes de transition sont affines. Nous
supposons que $X_i$ est quasi-compact et quasi-séparé pour tout $i \in I$.
Nous choisissons en outre un élément $0 \in I$.
\end{situation}

\begin{lemma}
\label{lemma-eventually-separated}
Dans la situation \ref{situation-descent}, supposons que
$\mathcal{X}_0 \to Y_0$ soit un morphisme d'un champ algébrique vers $Y_0$.
Supposons $\mathcal{X}_0$ quasi-compact et quasi-séparé. Si
$Y \times_{Y_0} \mathcal{X}_0 \to Y$ est séparé, alors
$Y_i \times_{Y_0} \mathcal{X}_0 \to Y_i$ est séparé pour tout $i \in I$
suffisamment grand.
\end{lemma}

\begin{proof}
Écrivons $\mathcal{X} = Y \times_{Y_0} \mathcal{X}_0$ et
$\mathcal{X}_i = Y_i \times_{Y_0} \mathcal{X}_0$. Choisissons un schéma
affine $U_0$ et un morphisme lisse surjectif
$U_0 \to \mathcal{X}_0$. Posons $U = Y \times_{Y_0} U_0$ et
$U_i = Y_i \times_{Y_0} U_0$. Alors $U$ et $U_i$ sont affines, tandis que
$U \to \mathcal{X}$ et $U_i \to \mathcal{X}_i$ sont lisses et surjectifs.
Posons $R_0 = U_0 \times_{\mathcal{X}_0} U_0$, puis
$R = Y \times_{Y_0} R_0$ et $R_i = Y_i \times_{Y_0} R_0$. Nous avons alors
$R = U \times_\mathcal{X} U$ et
$R_i = U_i \times_{\mathcal{X}_i} U_i$.

\medskip\noindent
Avec ces notations, remarquons que, si $\mathcal{X} \to Y$ est séparé, alors
$R \to U \times_Y U$ est propre : c'est le changement de base de
$\mathcal{X} \to \mathcal{X} \times_Y \mathcal{X}$ suivant
$U \times_Y U \to \mathcal{X} \times_Y \mathcal{X}$. Réciproquement,
$\mathcal{X}_i \to Y_i$ est séparé dès que
$R_i \to U_i \times_{Y_i} U_i$ est propre, car
$U_i \times_{Y_i} U_i \to \mathcal{X}_i \times_{Y_i} \mathcal{X}_i$ est
lisse et surjectif ; voir Propriétés des champs, lemme
\ref{stacks-properties-lemma-check-property-covering}. Remarquons que
$R_0 \to U_0 \times_{Y_0} U_0$ est localement de type fini et que $R_0$ est
quasi-compact et quasi-séparé. Par Limites d'espaces, lemme
\ref{spaces-limits-lemma-eventually-proper}, le morphisme
$R_i \to U_i \times_{Y_i} U_i$ est propre pour $i$ assez grand, ce qui achève
la démonstration.
\end{proof}







\section{Descente d'objets relatifs}
\label{section-descending-relative}

\noindent
Cette section est l'analogue de Limites d'espaces, section
\ref{spaces-limits-section-descending-relative}.

\begin{lemma}
\label{lemma-descend-a-stack-down}
Soit $I$ un ensemble filtrant. Soit $(X_i, f_{ii'})$ un système projectif
d'espaces algébriques indexé par $I$. Supposons que
\begin{enumerate}
\item les morphismes $f_{ii'} : X_i \to X_{i'}$ sont affines,
\item les espaces $X_i$ sont quasi-compacts et quasi-séparés.
\end{enumerate}
Soit $X = \lim X_i$. Si $\mathcal{X}$ est un champ algébrique de présentation
finie sur $X$, il existe un $i \in I$ et un champ algébrique $\mathcal{X}_i$
de présentation finie sur $X_i$ tels que
$\mathcal{X} \cong \mathcal{X}_i \times_{X_i} X$ comme champs algébriques sur
$X$.
\end{lemma}

\begin{proof}
Par Morphismes de champs, définition
\ref{stacks-morphisms-definition-locally-finite-presentation}, le morphisme
$\mathcal{X} \to X$ est quasi-compact, localement de présentation finie et
quasi-séparé. Comme $X$ et $\mathcal{X} \to X$ sont quasi-compacts, le champ
$\mathcal{X}$ est quasi-compact (Morphismes de champs, définition
\ref{stacks-morphisms-definition-quasi-compact}). Il existe donc un schéma
affine $U$ et un morphisme lisse surjectif $U \to \mathcal{X}$ (Propriétés
des champs, lemme \ref{stacks-properties-lemma-quasi-compact-stack}). Posons
$R = U \times_\mathcal{X} U$. Nous obtenons un groupoïde lisse en espaces
algébriques $(U, R, s, t, c)$ sur $X$ tel que $\mathcal{X} = [U/R]$ ; voir
Champs algébriques, lemme \ref{algebraic-lemma-stack-presentation}. Puisque
$\mathcal{X} \to X$ et $X$ sont quasi-séparés, $\mathcal{X}$ est
quasi-séparé (Morphismes de champs, lemme
\ref{stacks-morphisms-lemma-composition-separated}). Ainsi,
$R \to U \times U$ est quasi-compact et quasi-séparé (Morphismes de champs,
lemme \ref{stacks-morphisms-lemma-fibre-product-after-map}) ; par conséquent,
$R$ est un espace algébrique quasi-compact et quasi-séparé. D'autre part,
$U \to X$ est localement de présentation finie, et il en va donc de même de
$R \to X$ (car $s : R \to U$ est lisse, donc localement de présentation
finie). Le quintuplet $(U, R, s, t, c)$ est ainsi un objet groupoïde dans la
catégorie des espaces algébriques de présentation finie sur $X$. Par Limites
d'espaces, lemme \ref{spaces-limits-lemma-descend-finite-presentation}, il
existe un $i$ et un groupoïde en espaces algébriques
$(U_i, R_i, s_i, t_i, c_i)$ sur $X_i$ dont le changement de base à $X$ est
isomorphe à $(U, R, s, t, c)$. Quitte à augmenter $i$, nous pouvons supposer
$s_i$ et $t_i$ lisses ; voir Limites d'espaces, lemme
\ref{spaces-limits-lemma-descend-smooth}. Le champ quotient
$\mathcal{X}_i = [U_i/R_i]$ est algébrique (Champs algébriques, théorème
\ref{algebraic-theorem-smooth-groupoid-gives-algebraic-stack}).

\medskip\noindent
Il existe un morphisme $[U/R] \to [U_i/R_i]$ ; voir Groupoïdes en espaces,
lemme \ref{spaces-groupoids-lemma-quotient-stack-functorial}. Nous affirmons
que, joint aux morphismes $[U/R] \to X$ et $[U_i/R_i] \to X_i$ (Groupoïdes
en espaces, lemme \ref{spaces-groupoids-lemma-quotient-stack-arrows}), il
fournit un isomorphisme, c'est-à-dire une équivalence,
$$
[U/R] \longrightarrow [U_i/R_i] \times_{X_i} X
$$
Le morphisme correspondant
$$
[U/_{\!p}R] \longrightarrow [U_i/_{\!p}R_i] \times_{X_i} X
$$
au niveau des « préfaisceaux de groupoïdes » de Groupoïdes en espaces,
équation (\ref{spaces-groupoids-equation-quotient-stack}), est un
isomorphisme. L'affirmation résulte donc du fait que le passage au champ
associé commute aux produits fibrés ; voir Champs, lemme
\ref{stacks-lemma-stackification-fibre-product-fibred-categories}.
\end{proof}




\section[Immersion fermée et présentation finie]{Immersion fermée d'un morphisme de type fini dans un morphisme de présentation finie}
\label{section-finite-type-closed-in-finite-presentation}

\noindent
Cette section est l'analogue de Limites d'espaces, section
\ref{spaces-limits-section-finite-type-closed-in-finite-presentation}.

\begin{lemma}
\label{lemma-finite-type-closed-in-finite-presentation}
Soit $f : \mathcal{X} \to Y$ un morphisme d'un champ algébrique vers un espace
algébrique. Supposons que :
\begin{enumerate}
\item $f$ est de type fini et quasi-séparé ;
\item $Y$ est quasi-compact et quasi-séparé.
\end{enumerate}
Il existe alors un morphisme de présentation finie
$f' : \mathcal{X}' \to Y$ et une immersion fermée
$\mathcal{X} \to \mathcal{X}'$ de champs algébriques sur $Y$.
\end{lemma}

\begin{proof}
Écrivons $Y = \lim_{i \in I} Y_i$ comme limite projective d'espaces algébriques
indexés par un ensemble filtrant $I$, dont les morphismes de transition sont
affines et les $Y_i$ noethériens ; voir Limites d'espaces, proposition
\ref{spaces-limits-proposition-approximate}. Nous utiliserons les résultats de
Limites d'espaces, section
\ref{spaces-limits-section-finite-type-quasi-separated}.

\medskip\noindent
Choisissons une présentation $\mathcal{X} = [U/R]$. Notons
$(U, R, s, t, c, e, i)$ le groupoïde correspondant en espaces algébriques sur
$Y$. Nous pouvons et allons supposer $U$ affine. Alors $U$, $R$ et
$R \times_{s, U, t} R$ sont des espaces algébriques quasi-séparés et de type
fini sur $Y$. Nous avons deux morphismes $s, t : R \to U$, trois morphismes
$c : R \times_{s, U, t} R \to R$,
$\text{pr}_1 : R \times_{s, U, t} R \to R$,
$\text{pr}_2 : R \times_{s, U, t} R \to R$,
un morphisme $e : U \to R$ et, enfin, un morphisme $i : R \to R$. Ces
morphismes satisfont une liste d'axiomes explicitée dans Groupoïdes, section
\ref{groupoids-section-groupoids}.

\medskip\noindent
D'après Limites d'espaces, remarque
\ref{spaces-limits-remark-finite-type-gives-well-defined-system}
nous pouvons trouver un $i_0 \in I$ et des systèmes projectifs
\begin{enumerate}
\item $(U_i)_{i \geq i_0}$,
\item $(R_i)_{i \geq i_0}$,
\item $(T_i)_{i \geq i_0}$
\end{enumerate}
sur $(Y_i)_{i \geq i_0}$ tels que
$U = \lim_{i \geq i_0} U_i$,
$R = \lim_{i \geq i_0} R_i$ et
$R \times_{s, U, t} R = \lim_{i \geq i_0} T_i$
et tels qu'il existe des morphismes de systèmes projectifs
\begin{enumerate}
\item $(s_i)_{i \geq i_0} : (R_i)_{i \geq i_0} \to (U_i)_{i \geq i_0}$,
\item $(t_i)_{i \geq i_0} : (R_i)_{i \geq i_0} \to (U_i)_{i \geq i_0}$,
\item $(c_i)_{i \geq i_0} : (T_i)_{i \geq i_0} \to (R_i)_{i \geq i_0}$,
\item $(p_i)_{i \geq i_0} : (T_i)_{i \geq i_0} \to (R_i)_{i \geq i_0}$,
\item $(q_i)_{i \geq i_0} : (T_i)_{i \geq i_0} \to (R_i)_{i \geq i_0}$,
\item $(e_i)_{i \geq i_0} : (U_i)_{i \geq i_0} \to (R_i)_{i \geq i_0}$,
\item $(i_i)_{i \geq i_0} : (R_i)_{i \geq i_0} \to (R_i)_{i \geq i_0}$
\end{enumerate}
vérifiant
$s = \lim_{i \geq i_0} s_i$,
$t = \lim_{i \geq i_0} t_i$,
$c = \lim_{i \geq i_0} c_i$,
$\text{pr}_1 = \lim_{i \geq i_0} p_i$,
$\text{pr}_2 = \lim_{i \geq i_0} q_i$,
$e = \lim_{i \geq i_0} e_i$ et
$i = \lim_{i \geq i_0} i_i$.
Par Limites d'espaces, lemme
\ref{spaces-limits-lemma-morphism-good-diagram-smooth}
nous pouvons supposer $s_i$ et $t_i$ lisses, quitte à augmenter $i_0$.
Par Limites d'espaces, lemme
\ref{spaces-limits-lemma-morphism-good-diagram-flat}
nous pouvons supposer que les applications
$R \to U \times_{U_i, s_i} R_i$ données par $s$ et $R \to R_i$, ainsi que
$R \to U \times_{U_i, t_i} R_i$ données par $t$ et $R \to R_i$, sont des
isomorphismes pour tout $i \geq i_0$. Par Limites d'espaces, lemme
\ref{spaces-limits-lemma-good-diagram-fibre-product}, nous pouvons supposer les
diagrammes
$$
\xymatrix{
T_i \ar[r]_{q_i} \ar[d]_{p_i} & R_i \ar[d]^{t_i} \\
R_i \ar[r]^{s_i} & U_i
}
$$
cartésiens. L'assertion d'unicité de Limites d'espaces, lemme
\ref{spaces-limits-lemma-morphism-good-diagram}, garantit alors que, pour un
$i$ assez grand, les relations mentionnées plus haut entre les morphismes
$s, t, c, e, i$ sont satisfaites par $s_i, t_i, c_i, e_i, i_i$. Fixons un tel
$i$.

\medskip\noindent
Il s'ensuit que $(U_i, R_i, s_i, t_i, c_i, e_i, i_i)$ est un groupoïde lisse en
espaces algébriques sur $Y_i$. Ainsi, $\mathcal{X}_i = [U_i/R_i]$ est un champ
algébrique (Champs algébriques, théorème
\ref{algebraic-theorem-smooth-groupoid-gives-algebraic-stack}).
Le morphisme de groupoïdes
$$
(U, R, s, t, c, e, i) \to (U_i, R_i, s_i, t_i, c_i, e_i, i_i)
$$
au-dessus de $Y \to Y_i$ détermine un diagramme commutatif
$$
\xymatrix{
\mathcal{X} \ar[d] \ar[r] & \mathcal{X}_i \ar[d] \\
Y \ar[r] & Y_i
}
$$
(Groupoïdes en espaces, lemme
\ref{spaces-groupoids-lemma-quotient-stack-functorial}).
Nous affirmons que le morphisme $\mathcal{X} \to Y \times_{Y_i} \mathcal{X}_i$
est une immersion fermée. Cette affirmation achève la démonstration, car le
champ algébrique $\mathcal{X}_i \to Y_i$ est de présentation finie par
construction. Pour la prouver, remarquons que le diagramme de gauche
$$
\xymatrix{
U \ar[d] \ar[r] & U_i \ar[d] \\
\mathcal{X} \ar[r] & \mathcal{X}_i
}
\quad\quad
\xymatrix{
U \ar[d] \ar[r] & Y \times_{Y_i} U_i \ar[d] \\
\mathcal{X} \ar[r] & Y \times_{Y_i} \mathcal{X}_i
}
$$
est cartésien par Groupoïdes en espaces, lemme
\ref{spaces-groupoids-lemma-criterion-fibre-product}
et les résultats rappelés ci-dessus. Le diagramme commutatif de droite est donc
lui aussi cartésien. Le résultat voulu découle alors du fait que
$U \to Y \times_{Y_i} U_i$ est une immersion fermée, par construction du
système projectif $(U_i)$ dans Limites d'espaces, lemme
\ref{spaces-limits-lemma-limit-from-good-diagram}, du fait que
$Y \times_{Y_i} U_i \to Y \times_{Y_i} \mathcal{X}_i$ est lisse et surjectif,
et de Propriétés des champs, lemme
\ref{stacks-properties-lemma-check-immersion-covering}.
\end{proof}

\noindent
Il existe une variante pour les champs algébriques séparés.

\begin{lemma}
\label{lemma-separated-closed-in-finite-presentation}
Soit $f : \mathcal{X} \to Y$ un morphisme d'un champ algébrique vers un espace
algébrique. Supposons que :
\begin{enumerate}
\item $f$ est de type fini et séparé ;
\item $Y$ est quasi-compact et quasi-séparé.
\end{enumerate}
Il existe alors un morphisme séparé de présentation finie
$f' : \mathcal{X}' \to Y$ et une immersion fermée
$\mathcal{X} \to \mathcal{X}'$ de champs algébriques sur $Y$.
\end{lemma}

\begin{proof}
Nous appliquons d'abord exactement le même procédé que dans la démonstration du
lemme \ref{lemma-finite-type-closed-in-finite-presentation}, dont nous reprenons
les notations, pour construire l'immersion $\mathcal{X} \to \mathcal{X}'$ sous
la forme d'un morphisme
$\mathcal{X} \to \mathcal{X}' = Y \times_{Y_i} \mathcal{X}_i$, où
$\mathcal{X}_i = [U_i/R_i]$. Il suffit donc de montrer que
$\mathcal{X}_i \to Y_i$ est séparé pour $i$ assez grand. Autrement dit, il
suffit de montrer que
$\mathcal{X}_i \to \mathcal{X}_i \times_{Y_i} \mathcal{X}_i$ est propre pour
$i$ assez grand. Puisque le morphisme
$U_i \times_{Y_i} U_i \to \mathcal{X}_i \times_{Y_i} \mathcal{X}_i$ est
surjectif et lisse, et puisque
$R_i = \mathcal{X}_i
\times_{\mathcal{X}_i \times_{Y_i} \mathcal{X}_i} U_i \times_{Y_i} U_i$
il suffit de montrer que le morphisme
$(s_i, t_i) : R_i \to U_i \times_{Y_i} U_i$
est propre pour $i$ assez grand ; voir Propriétés des champs, lemme
\ref{stacks-properties-lemma-check-property-covering}.
Nous le démontrons au paragraphe suivant.

\medskip\noindent
Observons que $U \times_Y U \to Y$ est quasi-séparé et de type fini. Nous
pouvons donc appliquer la construction de Limites d'espaces, remarque
\ref{spaces-limits-remark-finite-type-gives-well-defined-system}
pour trouver un $i_1 \in I$ et un système projectif $(V_i)_{i \geq i_1}$ tels
que $U \times_Y U = \lim_{i \geq i_1} V_i$. Par Limites d'espaces, lemme
\ref{spaces-limits-lemma-good-diagram-fibre-product}, pour $i$ assez grand, la
fonctorialité de la construction appliquée aux projections
$U \times_Y U \to U$ donne des immersions fermées
$$
V_i \to U_i \times_{Y_i} U_i
$$
(Il existe ici un léger décalage : il faudrait en réalité remplacer $Y_i$ par
l'image schématique de $Y \to Y_i$, mais cela ne change manifestement pas le
produit fibré.) D'autre part, par Limites d'espaces, lemme
\ref{spaces-limits-lemma-morphism-good-diagram-proper}
la fonctorialité appliquée au morphisme propre
$(s, t) : R \to U \times_Y U$ (nous utilisons ici que $\mathcal{X}$ est séparé)
fournit des morphismes $R_i \to V_i$ propres pour $i$ assez grand. En les
composant, nous obtenons un morphisme propre
$R_i \to U_i \times_{Y_i} U_i$ pour tout $i$ assez grand. La fonctorialité de
la construction de Limites d'espaces, remarque
\ref{spaces-limits-remark-finite-type-gives-well-defined-system}
montre que ce morphisme coïncide avec $(s_i, t_i)$ pour $i$ assez grand, ce qui
achève la démonstration.
\end{proof}






\section{Morphismes universellement fermés}
\label{section-universally-closed}

\noindent
Cette section est l'analogue de Limites d'espaces, section
\ref{spaces-limits-section-universally-closed}.

\begin{lemma}
\label{lemma-separate}
Soit $g : Z \to Y$ un morphisme de schémas affines. Soit
$f : \mathcal{X} \to Y$ un morphisme quasi-compact de champs algébriques. Soit
$z \in Z$, et soit $T \subset |\mathcal{X} \times_Y Z|$ une partie fermée telle
que $z \not \in \Im(T \to |Z|)$. Si $\mathcal{X}$ est quasi-compact, il existe
un voisinage ouvert $V \subset Z$ de $z$, un diagramme commutatif
$$
\xymatrix{
V \ar[d] \ar[r]_a & Z' \ar[d]^b \\
Z \ar[r]^g & Y,
}
$$
et une partie fermée $T' \subset |X \times_Y Z'|$ tels que :
\begin{enumerate}
\item $Z'$ est un schéma affine de présentation finie sur $Y$ ;
\item en posant $z' = a(z)$, on a $z' \not \in \Im(T' \to |Z'|)$ ;
\item l'image réciproque de $T$ dans $|\mathcal{X} \times_Y V|$ s'envoie dans
$T'$ par $|\mathcal{X} \times_Y V| \to |\mathcal{X} \times_Y Z'|$.
\end{enumerate}
\end{lemma}

\begin{proof}
Nous allons déduire l'assertion du résultat correspondant pour les morphismes
de schémas. Comme $\mathcal{X}$ est quasi-compact, choisissons un schéma affine
$W$ et un morphisme lisse surjectif $W \to \mathcal{X}$. Soit
$T_W \subset |W \times_Y Z|$ l'image réciproque de $T$. Alors $z$ n'appartient
pas à l'image de $T_W$. Le cas des schémas (Limites, lemme
\ref{limits-lemma-separate}) fournit un voisinage ouvert $V \subset Z$ de $z$,
un diagramme commutatif de schémas
$$
\xymatrix{
V \ar[d] \ar[r]_a & Z' \ar[d]^b \\
Z \ar[r]^g & Y,
}
$$
et une partie fermée $T' \subset |W \times_Y Z'|$ tels que :
\begin{enumerate}
\item $Z'$ est un schéma affine de présentation finie sur $Y$ ;
\item en posant $z' = a(z)$, on a $z' \not \in \Im(T' \to |Z'|)$ ;
\item $T_1 = T_W \cap |W \times_Y V|$ s'envoie dans $T'$ par
$|W \times_Y V| \to |W \times_Y Z'|$.
\end{enumerate}
Le diagramme commutatif
$$
\xymatrix{
W \times_Y Z \ar[d] &
W \times_Y V \ar[l] \ar[rr]_{a_1} \ar[d]_c & &
W \times_Y Z' \ar[d]^q \\
\mathcal{X} \times_Y Z &
\mathcal{X} \times_Y V \ar[l] \ar[rr]^{a_2} & &
\mathcal{X} \times_Y Z'
}
$$
est formé de carrés cartésiens, et ses flèches verticales sont surjectives,
lisses et, a fortiori, ouvertes. Le carré de gauche montre que
$T_1 = T_W \cap |W \times_Y V|$ est l'image réciproque de
$T_2 = T \cap |\mathcal{X} \times_Y V|$ par $c$. Par Propriétés des champs, lemme
\ref{stacks-properties-lemma-points-cartesian}, on obtient
$a_1(T_1) = q^{-1}(a_2(T_2))$.
Par Topologie, lemme \ref{topology-lemma-open-morphism-quotient-topology}, on a
$$
q^{-1}\left(\overline{a_2(T_2)}\right) =
\overline{q^{-1}(a_2(T_2))} =
\overline{a_1(T_1)} \subset T'
$$
Comme $q$ est surjectif, l'image de $\overline{a_2(T_2)} \to |Z'|$ ne contient
pas $z'$, puisque cela vaut pour $T'$. Le diagramme ci-dessus, avec
$Z', V, a, b$, et la partie fermée
$\overline{a_2(T_2)} \subset |\mathcal{X} \times_Y Z'|$ répondent donc au
problème posé dans le lemme.
\end{proof}

\begin{lemma}
\label{lemma-test-universally-closed}
Soit $f : \mathcal{X} \to \mathcal{Y}$ un morphisme quasi-compact de champs
algébriques. Les conditions suivantes sont équivalentes :
\begin{enumerate}
\item $f$ est universellement fermé ;
\item pour tout morphisme $Z \to \mathcal{Y}$ localement de présentation finie,
où $Z$ est un schéma affine, l'application
$|\mathcal{X} \times_Y Z| \to |Z|$ est fermée ;
\item il existe un schéma $V$ et un morphisme lisse surjectif
$V \to \mathcal{Y}$ tels que l'application
$|\mathbf{A}^n \times (\mathcal{X} \times_\mathcal{Y} V)|
\to |\mathbf{A}^n \times V|$ soit fermée pour tout $n \geq 0$.
\end{enumerate}
\end{lemma}

\begin{proof}
Il est clair que (1) implique (2).

\medskip\noindent
Supposons (2). Choisissons un schéma $V$ qui soit réunion disjointe de schémas
affines, ainsi qu'un morphisme lisse surjectif $V \to \mathcal{Y}$. Pour montrer
que $f$ est universellement fermé, il suffit de montrer que le changement de
base $\mathcal{X} \times_\mathcal{Y} V \to V$ de $f$ est universellement fermé ;
voir Morphismes de champs, lemme
\ref{stacks-morphisms-lemma-universally-closed-local}.
Remarquons que la condition (2) vaut pour ce changement de base. Afin de prouver
que (2) implique (1), nous pouvons donc supposer que $Y = \mathcal{Y}$ est un
schéma affine.

\medskip\noindent
Supposons (2), et supposons que $\mathcal{Y} = Y$ est un schéma affine. Si $f$
n'est pas universellement fermé, il existe un schéma affine $Z$ sur $Y$ tel que
$|\mathcal{X} \times_Y Z| \to |Z|$ ne soit pas fermé ; voir Morphismes de
champs, lemme
\ref{stacks-morphisms-lemma-universally-closed-local}.
Il existe donc une partie fermée $T \subset |\mathcal{X} \times_Y Z|$ telle que
$\Im(T \to |Z|)$ ne soit pas fermée. Choisissons $z \in |Z|$ dans l'adhérence
de l'image de $T$, mais non dans cette image. Appliquons le lemme
\ref{lemma-separate}. Nous obtenons un voisinage ouvert $V \subset Z$, un
diagramme commutatif
$$
\xymatrix{
V \ar[d] \ar[r]_a & Z' \ar[d]^b \\
Z \ar[r]^g & Y,
}
$$
et une partie fermée $T' \subset |\mathcal{X} \times_Y Z'|$ tels que :
\begin{enumerate}
\item $Z'$ est un schéma affine de présentation finie sur $Y$ ;
\item en posant $z' = a(z)$, on a $z' \not \in \Im(T' \to |Z'|)$ ;
\item l'image réciproque de $T$ dans $|\mathcal{X} \times_Y V|$ s'envoie dans
$T'$ par
$|\mathcal{X} \times_Y V| \to |\mathcal{X} \times_Y Z'|$.
\end{enumerate}
Nous affirmons que $z'$ appartient à l'adhérence de $\Im(T' \to |Z'|)$. Il en
résulte que $|\mathcal{X} \times_Y Z'| \to |Z'|$ n'est pas fermée, contrairement
à (2). Autrement dit, cette affirmation montre que (2) implique (1). Pour la
vérifier, considérons le diagramme commutatif suivant :
$$
\xymatrix{
\mathcal{X} \times_Y Z \ar[d] &
\mathcal{X} \times_Y V \ar[l] \ar[d] \ar[r] &
\mathcal{X} \times_Y Z' \ar[d] \\
Z & V \ar[l] \ar[r]^a & Z'
}
$$
Soit $T_V \subset |\mathcal{X} \times_Y V|$ l'image réciproque de $T$. Par
Propriétés des champs, lemme \ref{stacks-properties-lemma-points-cartesian},
l'image de $T_V$ dans $|V|$ est l'image réciproque de l'image de $T$ dans
$|Z|$. Puisque $z$ appartient à l'adhérence de l'image de $T \to |Z|$ et que
$|V| \to |Z|$ est ouverte, $z$ appartient donc à l'adhérence de l'image de
$T_V \to |V|$. Comme l'image de $T_V$ dans
$|\mathcal{X} \times_Y Z'|$ est contenue dans $|T'|$, il s'ensuit immédiatement
que $z' = a(z)$ appartient à l'adhérence de l'image de $T'$.

\medskip\noindent
Il est clair que (1) implique (3). Soit $V \to \mathcal{Y}$ comme dans (3). Si
nous montrons que $\mathcal{X} \times_Y V \to V$ est universellement fermé,
alors $f$ est universellement fermé par Morphismes de champs, lemme
\ref{stacks-morphisms-lemma-universally-closed-local}.
Il suffit donc de montrer que $f : \mathcal{X} \to \mathcal{Y}$ satisfait (2)
lorsque $f$ est un morphisme quasi-compact de champs algébriques, que
$\mathcal{Y} = Y$ est un schéma et que
$|\mathbf{A}^n \times \mathcal{X}| \to |\mathbf{A}^n \times Y|$
soit fermée pour tout $n$. Soit $Z \to Y$ localement de présentation finie,
où $Z$ est un schéma affine. Nous devons montrer que l'application
$|\mathcal{X} \times_Y Z| \to |Z|$ est fermée. Comme $Y$ est un schéma, que
$Z$ est affine et que $Z \to Y$ est localement de présentation finie, il
existe une immersion $Z \to \mathbf{A}^n \times Y$ ; voir Morphismes, lemme
\ref{morphisms-lemma-quasi-affine-finite-type-over-S}. Considérons le diagramme
cartésien
$$
\vcenter{
\xymatrix{
\mathcal{X} \times_Y Z \ar[d] \ar[r] &
\mathbf{A}^n \times \mathcal{X} \ar[d] \\
Z \ar[r] & \mathbf{A}^n \times Y
}
}
\quad
\begin{matrix}
\text{qui induit le} \\
\text{carré cartésien}
\end{matrix}
\quad
\vcenter{
\xymatrix{
|\mathcal{X} \times_Y Z| \ar[d] \ar[r] &
|\mathbf{A}^n \times \mathcal{X}| \ar[d] \\
|Z| \ar[r] & |\mathbf{A}^n \times Y|
}
}
$$
des espaces topologiques, dont les flèches horizontales sont des
homéomorphismes sur des parties localement fermées (Propriétés des champs, lemme
\ref{stacks-properties-lemma-subspace-induced-topology}).
Toute partie fermée $T$ de $|X \times_Y Z|$ est donc l'image réciproque d'une
partie fermée $T'$ de $|\mathbf{A}^n \times Y|$. Comme, par hypothèse, l'image
de $T'$ dans $|\mathbf{A}^n \times X|$ est fermée, nous en concluons que l'image
de $T$ dans $|Z|$ est fermée, comme voulu.
\end{proof}








\begin{multicols}{2}[\section{Autres chapitres}]
\noindent
Préliminaires
\begin{enumerate}
\item \hyperref[introduction-section-phantom]{Introduction}
\item \hyperref[conventions-section-phantom]{Conventions}
\item \hyperref[sets-section-phantom]{Théorie des ensembles}
\item \hyperref[categories-section-phantom]{Catégories}
\item \hyperref[topology-section-phantom]{Topologie}
\item \hyperref[sheaves-section-phantom]{Faisceaux sur les espaces}
\item \hyperref[sites-section-phantom]{Sites et faisceaux}
\item \hyperref[stacks-section-phantom]{Champs}
\item \hyperref[fields-section-phantom]{Corps}
\item \hyperref[algebra-section-phantom]{Algèbre commutative}
\item \hyperref[brauer-section-phantom]{Groupes de Brauer}
\item \hyperref[homology-section-phantom]{Algèbre homologique}
\item \hyperref[derived-section-phantom]{Catégories dérivées}
\item \hyperref[simplicial-section-phantom]{Méthodes simpliciales}
\item \hyperref[more-algebra-section-phantom]{Compléments d'algèbre}
\item \hyperref[smoothing-section-phantom]{Lissage des morphismes d'anneaux}
\item \hyperref[modules-section-phantom]{Faisceaux de modules}
\item \hyperref[sites-modules-section-phantom]{Modules sur les sites}
\item \hyperref[injectives-section-phantom]{Objets injectifs}
\item \hyperref[cohomology-section-phantom]{Cohomologie des faisceaux}
\item \hyperref[sites-cohomology-section-phantom]{Cohomologie sur les sites}
\item \hyperref[dga-section-phantom]{Algèbre différentielle graduée}
\item \hyperref[dpa-section-phantom]{Algèbre à puissances divisées}
\item \hyperref[sdga-section-phantom]{Faisceaux différentiels gradués}
\item \hyperref[hypercovering-section-phantom]{Hyperrecouvrements}
\end{enumerate}
Schémas
\begin{enumerate}
\setcounter{enumi}{25}
\item \hyperref[schemes-section-phantom]{Schémas}
\item \hyperref[constructions-section-phantom]{Constructions de schémas}
\item \hyperref[properties-section-phantom]{Propriétés des schémas}
\item \hyperref[morphisms-section-phantom]{Morphismes de schémas}
\item \hyperref[coherent-section-phantom]{Cohomologie des schémas}
\item \hyperref[divisors-section-phantom]{Diviseurs}
\item \hyperref[limits-section-phantom]{Limites de schémas}
\item \hyperref[varieties-section-phantom]{Variétés}
\item \hyperref[topologies-section-phantom]{Topologies sur les schémas}
\item \hyperref[descent-section-phantom]{Descente}
\item \hyperref[perfect-section-phantom]{Catégories dérivées de schémas}
\item \hyperref[more-morphisms-section-phantom]{Compléments sur les morphismes}
\item \hyperref[flat-section-phantom]{Compléments sur la platitude}
\item \hyperref[groupoids-section-phantom]{Schémas en groupoïdes}
\item \hyperref[more-groupoids-section-phantom]{Compléments sur les schémas en groupoïdes}
\item \hyperref[etale-section-phantom]{Morphismes étales de schémas}
\end{enumerate}
Thèmes de théorie des schémas
\begin{enumerate}
\setcounter{enumi}{41}
\item \hyperref[chow-section-phantom]{Homologie de Chow}
\item \hyperref[intersection-section-phantom]{Théorie de l'intersection}
\item \hyperref[pic-section-phantom]{Schémas de Picard des courbes}
\item \hyperref[weil-section-phantom]{Théories cohomologiques de Weil}
\item \hyperref[adequate-section-phantom]{Modules adéquats}
\item \hyperref[dualizing-section-phantom]{Complexes dualisants}
\item \hyperref[duality-section-phantom]{Dualité pour les schémas}
\item \hyperref[discriminant-section-phantom]{Discriminants et différentes}
\item \hyperref[derham-section-phantom]{Cohomologie de de Rham}
\item \hyperref[local-cohomology-section-phantom]{Cohomologie locale}
\item \hyperref[algebraization-section-phantom]{Géométrie algébrique et formelle}
\item \hyperref[curves-section-phantom]{Courbes algébriques}
\item \hyperref[resolve-section-phantom]{Résolution des surfaces}
\item \hyperref[models-section-phantom]{Réduction semi-stable}
\item \hyperref[functors-section-phantom]{Foncteurs et morphismes}
\item \hyperref[equiv-section-phantom]{Catégories dérivées de variétés}
\item \hyperref[pione-section-phantom]{Groupes fondamentaux de schémas}
\item \hyperref[etale-cohomology-section-phantom]{Cohomologie étale}
\item \hyperref[crystalline-section-phantom]{Cohomologie cristalline}
\item \hyperref[proetale-section-phantom]{Cohomologie pro-étale}
\item \hyperref[relative-cycles-section-phantom]{Cycles relatifs}
\item \hyperref[more-etale-section-phantom]{Compléments de cohomologie étale}
\item \hyperref[trace-section-phantom]{La formule des traces}
\end{enumerate}
Espaces algébriques
\begin{enumerate}
\setcounter{enumi}{64}
\item \hyperref[spaces-section-phantom]{Espaces algébriques}
\item \hyperref[spaces-properties-section-phantom]{Propriétés des espaces algébriques}
\item \hyperref[spaces-morphisms-section-phantom]{Morphismes d'espaces algébriques}
\item \hyperref[decent-spaces-section-phantom]{Espaces algébriques décents}
\item \hyperref[spaces-cohomology-section-phantom]{Cohomologie des espaces algébriques}
\item \hyperref[spaces-limits-section-phantom]{Limites d'espaces algébriques}
\item \hyperref[spaces-divisors-section-phantom]{Diviseurs sur les espaces algébriques}
\item \hyperref[spaces-over-fields-section-phantom]{Espaces algébriques sur des corps}
\item \hyperref[spaces-topologies-section-phantom]{Topologies sur les espaces algébriques}
\item \hyperref[spaces-descent-section-phantom]{Descente et espaces algébriques}
\item \hyperref[spaces-perfect-section-phantom]{Catégories dérivées d'espaces}
\item \hyperref[spaces-more-morphisms-section-phantom]{Compléments sur les morphismes d'espaces}
\item \hyperref[spaces-flat-section-phantom]{Platitude sur les espaces algébriques}
\item \hyperref[spaces-groupoids-section-phantom]{Groupoïdes dans les espaces algébriques}
\item \hyperref[spaces-more-groupoids-section-phantom]{Compléments sur les groupoïdes dans les espaces}
\item \hyperref[bootstrap-section-phantom]{Amorçage}
\item \hyperref[spaces-pushouts-section-phantom]{Sommes amalgamées d'espaces algébriques}
\end{enumerate}
Thèmes de géométrie
\begin{enumerate}
\setcounter{enumi}{81}
\item \hyperref[spaces-chow-section-phantom]{Groupes de Chow des espaces}
\item \hyperref[groupoids-quotients-section-phantom]{Quotients de groupoïdes}
\item \hyperref[spaces-more-cohomology-section-phantom]{Compléments sur la cohomologie des espaces}
\item \hyperref[spaces-simplicial-section-phantom]{Espaces simpliciaux}
\item \hyperref[spaces-duality-section-phantom]{Dualité pour les espaces}
\item \hyperref[formal-spaces-section-phantom]{Espaces algébriques formels}
\item \hyperref[restricted-section-phantom]{Algébrisation des espaces formels}
\item \hyperref[spaces-resolve-section-phantom]{Résolution des surfaces revisitée}
\end{enumerate}
Théorie des déformations
\begin{enumerate}
\setcounter{enumi}{89}
\item \hyperref[formal-defos-section-phantom]{Théorie formelle des déformations}
\item \hyperref[defos-section-phantom]{Théorie des déformations}
\item \hyperref[cotangent-section-phantom]{Le complexe cotangent}
\item \hyperref[examples-defos-section-phantom]{Problèmes de déformation}
\end{enumerate}
Champs algébriques
\begin{enumerate}
\setcounter{enumi}{93}
\item \hyperref[algebraic-section-phantom]{Champs algébriques}
\item \hyperref[examples-stacks-section-phantom]{Exemples de champs}
\item \hyperref[stacks-sheaves-section-phantom]{Faisceaux sur les champs algébriques}
\item \hyperref[criteria-section-phantom]{Critères de représentabilité}
\item \hyperref[artin-section-phantom]{Axiomes d'Artin}
\item \hyperref[quot-section-phantom]{Espaces de Quot et de Hilbert}
\item \hyperref[stacks-properties-section-phantom]{Propriétés des champs algébriques}
\item \hyperref[stacks-morphisms-section-phantom]{Morphismes de champs algébriques}
\item \hyperref[stacks-limits-section-phantom]{Limites de champs algébriques}
\item \hyperref[stacks-cohomology-section-phantom]{Cohomologie des champs algébriques}
\item \hyperref[stacks-perfect-section-phantom]{Catégories dérivées de champs}
\item \hyperref[stacks-introduction-section-phantom]{Introduction aux champs algébriques}
\item \hyperref[stacks-more-morphisms-section-phantom]{Compléments sur les morphismes de champs}
\item \hyperref[stacks-geometry-section-phantom]{La géométrie des champs}
\end{enumerate}
Thèmes de théorie des espaces de modules
\begin{enumerate}
\setcounter{enumi}{107}
\item \hyperref[moduli-section-phantom]{Champs de modules}
\item \hyperref[moduli-curves-section-phantom]{Espaces de modules de courbes}
\end{enumerate}
Divers
\begin{enumerate}
\setcounter{enumi}{109}
\item \hyperref[examples-section-phantom]{Exemples}
\item \hyperref[exercises-section-phantom]{Exercices}
\item \hyperref[guide-section-phantom]{Guide bibliographique}
\item \hyperref[desirables-section-phantom]{Desiderata}
\item \hyperref[coding-section-phantom]{Style de programmation}
\item \hyperref[obsolete-section-phantom]{Obsolète}
\item \hyperref[fdl-section-phantom]{Licence de documentation libre GNU}
\item \hyperref[index-section-phantom]{Index généré automatiquement}
\end{enumerate}
\end{multicols}


\bibliography{my}
\bibliographystyle{amsalpha}

\end{document}
